\chapter{First arrival, spectral packets, algebraic rarity, and tensor regularization}\label{ch:spectral-algebraic}
\index{Green kernel}
\index{first-arrival valuation}
\index{word map}
\index{algebraic defect locus}
\index{border rank}
\index{tensor capacity}


\section{Word distance as first arrival}
Let $G$ be finite and let $S=S^{-1}$ be a symmetric generating multiset.  In the right regular representation, define
\[
A_S=\sum_{s\in S}\rho(s).
\]
Then
\[
\langle\delta_g,A_S^k\delta_e\rangle
\]
counts length-$k$ words ending at $g$.

\begin{theorem}[Endpoint Green valuation]
The word distance is
\[
d_S(e,g)
=
\ord_{z=0}
\left\langle
\delta_g,(I-zA_S)^{-1}\delta_e
\right\rangle.
\]
\end{theorem}
\begin{proof}
Expand the resolvent as a formal power series:
\[
(I-zA_S)^{-1}=\sum_{k\ge0}z^kA_S^k.
\]
The first nonzero endpoint coefficient occurs at the shortest word length.
\end{proof}

For strictly positive integer costs, replace each generator $s$ by the monomial $z^{c(s)}\rho(s)$.  Every coefficient then counts finitely many words of a given total cost, and the same formal inverse is defined.  Zero-cost cycles require separate treatment: they can create infinitely many words of cost zero, so the ordinary counting series need not exist.  The finite atlases with free gates use shortest-path search with zero-cost edges instead.  Multivariate variables can track several resources when the grading gives every generator positive total degree.

\section{Why eigenvalues alone are insufficient}
If
\[
A_S=\sum_\lambda\lambda\Pi_\lambda,
\]
then
\[
\mathcal G_{S,g}(z)
=
\sum_\lambda
\frac{\langle\delta_g,\Pi_\lambda\delta_e\rangle}{1-\lambda z}.
\]
The eigenvalues and multiplicities do not determine the endpoint coefficients
\[
\langle\delta_g,\Pi_\lambda\delta_e\rangle.
\]
Two Cayley or Schreier systems can therefore be cospectral and still have different first-arrival behavior at distinguished targets.

\begin{principle}[Embedded packet principle]
For threshold questions, a spectral label is usually weaker than an embedded spectral projector.  The matrix coefficients of the projector relative to the initial and target states retain positional information that multiplicities erase.
\end{principle}

This distinction recurs in the Latin-square example discussed later.  The natural commuting row, column, and symbol operators have universal eigenvalue packets at fixed order, while the embedded incidence and transition data still vary.

\section{Subgroup-holonomy Fourier packets}
Let $H\le G$.  Let $\mu_k(h)$ count ambient words of length $k$ with product $h\in H$.  For every irreducible representation $\tau\in\widehat H$, define
\[
\widehat\mu_k(\tau)
=
\sum_{h\in H}\mu_k(h)\tau(h).
\]
Fourier inversion gives
\[
\mu_k(h)
=
\frac1{|H|}
\sum_{\tau\in\widehat H}
\dim(\tau)
\operatorname{tr}
\bigl(
\tau(h^{-1})\widehat\mu_k(\tau)
\bigr).
\]
The matrix-valued Green packet
\[
\widehat{\mathcal G}_\tau(z)
=
\sum_{k\ge0}\widehat\mu_k(\tau)z^k
\]
recovers every endpoint series after Fourier inversion.

The trivial representation records only the total number of loops returning to the subgroup.  Nontrivial $H$-types retain their holonomy.  This yields a hierarchy:
\[
\text{spherical return data}
\prec
\text{type-resolved holonomy packets}
\prec
\text{exact endpoint realization}.
\]

\section{Word maps and product loci}
Let $K$ be a field and let $X\subseteq G$ be an algebraic gate family.  For a fixed length $L$, multiplication defines
\[
\mu_L:X^L\to G,
\qquad
(x_L,\ldots,x_1)\mapsto x_L\cdots x_1.
\]
The exact length-$L$ locus is $\mu_L(X^L)$, and the length-at-most-$L$ locus is a finite union of such images when the gate family has finitely many components.

\begin{theorem}[Dimension per gate]
Suppose every irreducible component of $X$ has dimension at most $r$.  Then the Zariski closure of the locus reachable in at most $L$ gates has dimension at most $rL$.
\end{theorem}
\begin{proof}
Every component of $X^L$ has dimension at most $rL$.  A morphism cannot increase dimension.  A finite union has dimension equal to the maximum dimension of its components.
\end{proof}

For scaled transvections, $r=1$.  This immediately yields generic lower bounds in a target variety whose dimension exceeds $L$.

\section{Generic rigidity}
Work over an algebraically closed field.  Suppose the restricted target variety $V$ is irreducible of dimension $N$, every target has a restricted realization using at most $N$ unit-cost gates, and the ambient alphabet is an algebraic family with finitely many components of dimension at most one.  Then ambient realizations of length below $N$ occupy a proper closed subset of $V$.

\begin{theorem}[Abstract generic rigidity]
Under the preceding hypotheses, there is a nonempty Zariski-open subset $\Omega\subseteq V$ such that
\[
\ellRes(U)=\ellAmb(U)=N
\qquad(U\in\Omega).
\]
\end{theorem}
The unitriangular theorem in the core chapter is a concrete instance with
\[
N=\binom n2.
\]

\section{Defect codimension}
If every target has restricted cost at most $N$, then defect at least $d$ forces ambient cost at most $N-d$.  The dimension-per-gate theorem gives
\[
\codim_V\overline{\{U:\Delta(U)\ge d\}}
\ge d
\]
when each ambient gate carries one parameter.

\begin{principle}[Rare but amplifiable]
At fixed dimension, large defects require increasingly special algebraic conditions.  Under direct sums, recursive lifts, or tensor products, a point on a thin exceptional stratum can nevertheless seed a macroscopic family.
\end{principle}

This is one of the central tensions of the subject.  Catalysis is neither generic nor negligible.

\section{Finite-field sparsity}
A word-counting analogue applies over $\F_q$.  If the ambient alphabet has at most $M(q-1)$ nontrivial scaled gates, then the number of targets reached in $L$ gates is at most
\[
\sum_{j=0}^L[M(q-1)]^j.
\]
When the target space has $q^N$ points and defect $d$ implies ambient length at most $N-d$, the crude bound is
\[
\Prob[\Delta\ge d]
=O_{N,M}(q^{-d}).
\]
The estimate can be sharpened by orbit counting, stabilizers, and the geometry of word-map fibers.

\section{Direct sums and monoidal amplification}
Suppose targets admit a direct sum $U\oplus V$ and a homomorphism from the large restricted group onto the two factor groups sends $U\oplus V$ to $(U,V)$.  Require the aggregate generator bound
\[
\ellRes(\pi_1(s))+\ellRes(\pi_2(s))\le c(s)
\]
for every large restricted generator $s$.  Assume that both restricted and ambient factor circuits embed into the corresponding large model without increasing cost.  Then
\[
\ellRes(U\oplus V)
=
\ellRes(U)+\ellRes(V).
\]
Indeed, projecting any word and summing the two factor bounds gives the lower bound, while concatenating the embedded factor circuits gives the upper bound.  Separate one-Lipschitz projections alone do not imply this sum inequality.
Under the ambient embedding assumption, the constructive upper bound is
\[
\ellAmb(U\oplus V)
\le
\ellAmb(U)+\ellAmb(V),
\]
so
\[
\Delta(U\oplus V)
\ge
\Delta(U)+\Delta(V).
\]
One finite defect therefore gives an unbounded additive family.

\section{Tensor products and normalized rates}
Let $U$ have dimension $m$ and $V$ dimension $n$.  In linear transvection models with tree-tight certificates,
\[
\ellRes(U\otimes V)
=
n\ellRes(U)+m\ellRes(V).
\]
The ambient model satisfies the corresponding inequality.  Normalizing by dimension gives an entropy-like law:
\[
\frac{\ellRes(U\otimes V)}{mn}
=
\frac{\ellRes(U)}m+
\frac{\ellRes(V)}n.
\]
The normalized defect is superadditive under tensor products.

\begin{definition}[Tensor synergy]
Define
\[
\Syn(U,V)
=
n\ellAmb(U)+m\ellAmb(V)-\ellAmb(U\otimes V).
\]
Positive synergy means that the joint ambient realization improves on independent lifted circuits.
\end{definition}
Restricted tree-tight metrics have zero tensor synergy.  Ambient synergy is a possible source of collective activation.

\begin{definition}[Stable tensor cost]
For a size-$m$ target $U$, define
\[
\tau_\otimes(U)
=
\lim_{r\to\infty}
\frac{\ellAmb(U^{\otimes r})}{r m^{r-1}},
\]
when the subadditivity argument applies.  Define
\[
\kappa_\otimes(U)
=
\ellRes(U)-\tau_\otimes(U).
\]
\end{definition}
Fekete's lemma gives the limit from the tensor subadditivity inequality.

\begin{figure}[ht]
\centering
\includegraphics[width=.72\textwidth]{tensor_capacity.pdf}
\caption{Schematic stable tensor capacity.  The figure illustrates definitions, not measured data.}
\end{figure}

\begin{definition}[Tensor activation]
A target is tensor activated if
\[
\Delta(U)=0
\quad\text{but}\quad
\kappa_\otimes(U)>0.
\]
\end{definition}
This is analogous to regularized capacities and asymptotic ranks.  It asks whether collective synthesis can unlock a catalyst invisible in one copy.

\section{Border catalysis}
Let $X_{\mathrm{res}}$ and $X_{\mathrm{amb}}$ be algebraic gate varieties.  The exact product rank is
\[
\operatorname{rk}_{X}(U)
=
\min\{L:U\in X^L\},
\]
while border product rank is
\[
\underline{\operatorname{rk}}_{X}(U)
=
\min\{L:U\in\overline{X^L}\}.
\]
The four possible comparisons are:
\begin{enumerate}
\item exact and border defects both positive,
\item exact defect positive but border defect zero,
\item one-copy defect zero but tensor-regularized defect positive,
\item ambient and restricted border ranks both smaller than exact ranks but by different amounts.
\end{enumerate}
Each requires a separate certificate.  Numerical approximation alone does not distinguish them.

\section{A packet-resolved threshold principle}
Combining the preceding constructions gives the following architecture.
\begin{enumerate}
\item A coarse commutative algebra supplies spherical spectral packets.
\item Embedded projectors retain incidence and endpoint position.
\item Noncommutative transition data record holonomy.
\item A realization theorem decides whether the packet has a restricted representative and at what cost.
\item Direct sums and tensor products propagate local defects.
\end{enumerate}

\begin{figure}[ht]
\centering
\includegraphics[width=.86\textwidth]{packet_hierarchy.pdf}
\caption{The spectral-to-realization hierarchy.}
\label{fig:packet-hierarchy}
\end{figure}

\begin{boundarybox}
A packet hierarchy organizes evidence.  It does not replace the final lower bound, positivity argument, orthogonality construction, or descent theorem.  Classification and realization remain distinct layers.
\end{boundarybox}
